\documentclass{article}
\usepackage[utf8]{inputenc}
\usepackage[T1]{fontenc}
\usepackage{geometry}
\geometry{
    a4paper,
    left=2.5cm,
    right=2.5cm,
    top=2.5cm,
    bottom=2.5cm
}
\usepackage{graphicx}
\usepackage{float}
\usepackage{hyperref}
\usepackage[english]{babel}
\usepackage{amsmath}
\usepackage{amssymb}
\usepackage{array}
\usepackage{booktabs}
\usepackage{caption}
\usepackage{subcaption}
\usepackage{enumitem}
\usepackage{fancyhdr}
\usepackage{titlesec}
\usepackage{url}
\usepackage{textcomp}
\usepackage{xcolor}
\usepackage{listings}
\usepackage{multirow}
\usepackage{longtable}
\usepackage{tabularx}
\usepackage{times}

% Set up page style
\pagestyle{fancy}
\fancyhf{}
\rfoot{\thepage}
\renewcommand{\headrulewidth}{0pt}

% Title formatting
\titleformat{\chapter}[display]{\normalfont\bfseries}{\chaptertitle\thechapter}{0pt}{\Huge}
\titlespacing*{\chapter}{0pt}{50pt}{40pt}

% Document info
\title{Template Document}
\author{Academic Support Agent}
\date{\today}

\begin{document}
\maketitle

% --------------------------------------------------------------
%                         Start here
% --------------------------------------------------------------
 
\title{Weekly Homework 1}
\author{Tony Stark\\ %replace with your name
Foundations of Mathematics}

\maketitle

\noindent For this assignment, I watched the following videos:

\begin{enumerate} %the first video is required (but you don't have to watch it first), you can choose the other 4 from the list on the Course Materials page
\item Grit: the power of passion and perseverance by Angela Lee Duckworth  
\item Replace with title for video 2
\item Replace with title for video 3
\item Replace with title for video 4
\item Replace with title for video 5
\end{enumerate}

Replace this text with your reflection.  Your total reflection should be at least 15 sentences long. You should comment on each of the videos you watched, but rather than reflecting on each video separately, try to reflect on growth mindset, productive failure, and grit, in general.
 
% --------------------------------------------------------------
%     You don't have to mess with anything below this line.
% --------------------------------------------------------------
\end{document}